\documentclass[11pt]{article}

\usepackage[utf8]{inputenc}
\usepackage[T1]{fontenc}
\usepackage{lmodern}
\usepackage{geometry}
\usepackage{amsmath,amssymb,amsfonts,amsthm,mathtools}
\usepackage{enumitem}
\usepackage{microtype}
\usepackage{hyperref}

\geometry{margin=1in}
\hypersetup{colorlinks=true, linkcolor=black, urlcolor=blue, citecolor=black}
\setlist[enumerate]{topsep=0.35em,itemsep=0.3em}
\setlist[itemize]{topsep=0.3em,itemsep=0.25em}

\newcommand{\Aether}{\AE{}ther}
\newcommand{\Aflow}{\AE{}ther-flow}
\newcommand{\TheoryProgram}{\Aether{} / \Aflow{} framework}
\newcommand{\Sub}{\mathcal{A}}
\newcommand{\PhiFlow}{\Phi_{\lambda}}
\newcommand{\Data}{\mathcal{D}_{\mathrm{loc}}}
\newcommand{\Map}{\mathcal{L}_{\mathrm{loc}}}
\newcommand{\Constraints}{\mathcal{C}_{\mathrm{loc}}}
\newcommand{\Front}{\mathcal{F}}
\newcommand{\Order}{\mathcal{S}}
\newcommand{\Obs}{\mathcal{O}}

\newtheorem{auditfinding}{Audit finding}
\newtheorem{blocker}{Promotion blocker}
\newtheorem{repairburden}{Repair burden}
\newtheorem{auditverdict}{Audit verdict}

\title{\textbf{Localization Smuggling Audit}\\[0.35em]
\large Target-Import Review of the Gate 0 Localization Candidate Packet}
\author{Smuggling Auditor AgentJob AJ-RT-20260608-004-001}
\date{June 8, 2026}

\begin{document}

\maketitle

\begin{abstract}
This draft audits the localization candidate packet for hidden target imports.
The audit focuses on the source front relation \(\Front\), the observer
coincidence response map \(\kappa_{\Gamma}\), the clock-scale calibration
\(\alpha_{\Gamma}\), and observer-chain selection. The conclusion is
conservative. The packet is not shown to be target-smuggled as written, because
it explicitly labels its structures as source-side and conditional. However,
it is not promotion-ready. Its current definitions leave open three
promotion-blocking routes by which the target null cone, target clock
normalization, or target observer protocol could be imported under source-side
names.
\end{abstract}

\section{Scope and Audit Standard}

This artifact is a draft/control output for task \texttt{RT-20260608-004}. It
audits
\[
\Data
=
(\Sub,\PhiFlow,\Gamma,\Order_{\Gamma},\Front,\kappa_{\Gamma},
\alpha_{\Gamma},\Map,\Constraints)
\]
from the localization candidate packet. It does not edit canonical ontology
TeX, does not promote a derivation, and does not issue a Gate Chair verdict.

The audit standard is source discipline. A structure passes only if it can be
specified from source objects, source order, source dynamics, or source
operational protocols without assuming the target Lorentzian metric, target
null cones, target proper time, target observer congruences, or the
Einsteinian field equations as primitives.

\section{Primary Findings}

\begin{auditfinding}[The front relation is named but not generated]
The packet states that \(\Front\subset\Sub\times\Sub\) records first influence
under \(\PhiFlow\), and it correctly says that \(\Front\) is not a metric null
cone. This is a necessary anti-smuggling statement, but not a sufficient
source definition. The current packet does not yet give a source law that
computes or constrains \(\Front\) from \((\Sub,\PhiFlow,\Order)\).
\end{auditfinding}

The risk is not the symbol \(\Front\). The risk is a later choice of \(\Front\)
that is tuned to produce the benchmark light cone. If the front relation is
selected because its localized image is already a Lorentzian null cone, then
the target causal structure has been imported rather than derived.

\begin{blocker}[Front-generation blocker]
Before promotion, the project must state a source-only generation rule for
\(\Front\). The rule must be expressible before comparison with any target
metric. It may then be tested for smooth double-front behavior, but the double
front cannot be imposed as the defining axiom.
\end{blocker}

\begin{auditfinding}[The coincidence map risks importing radar coordinates]
The packet defines
\[
\kappa_{\Gamma}:U_{\Gamma}\rightarrow
(\Order_{\Gamma},r^1_{\Gamma},r^2_{\Gamma},r^3_{\Gamma})
\]
by repeated interactions between an observer chain and nearby source fronts.
This is directionally correct: it tries to build labels operationally rather
than by a spatial metric. The gap is that the interaction protocol is not yet
specified. Without a source-only protocol, \(\kappa_{\Gamma}\) can silently
become a radar-coordinate construction that presupposes relativistic signal
conventions.
\end{auditfinding}

The danger point is the word ``nearby.'' Nearness is harmless if it is an
interaction domain defined by source adjacency, response amplitude, or
substrate transition accessibility. It is smuggling if it depends on a target
spatial distance, simultaneity slice, or preexisting light-signal convention.

\begin{blocker}[Coincidence-protocol blocker]
Before promotion, the candidate needs a source-only coincidence protocol for
\(\kappa_{\Gamma}\). The protocol must define response labels from observable
source interactions and must state its invariances under relabeling of
\(\Sub\), \(\lambda\), and the response coordinates.
\end{blocker}

\begin{auditfinding}[Clock scale is underdetermined]
The packet correctly separates order \(\Order_{\Gamma}\) from target proper
time and assigns scale responsibility to \(\alpha_{\Gamma}\). This avoids an
immediate clock-smuggling error. The unresolved risk is that
\(\alpha_{\Gamma}\) is only described as fixed by stable source cycles. The
packet does not yet define which cycles qualify, how their units compose, or
why the calibration is independent of target proper-time normalization.
\end{auditfinding}

A source cycle can calibrate a clock only after the model supplies a rule for
cycle identity, repeatability, stability, and comparison across observer
chains. If those rules are chosen to reproduce the benchmark clock after the
fact, the target clock has been imported.

\begin{blocker}[Clock-scale blocker]
Before promotion, the candidate must define \(\alpha_{\Gamma}\) by a
source-only cycle law. The law must state dimensions and must survive arbitrary
relabeling of the source order. It may be compared with target proper time only
after it is independently specified.
\end{blocker}

\begin{auditfinding}[The smooth double-front test is admissible only as a test]
The candidate asks whether the localized image of \(\Front\) becomes a smooth
double-front structure with one invariant signal speed. This is an admissible
failure test. It is not admissible as a source axiom. If smooth Lorentzian cone
structure is required at definition time, the test collapses into target
import.
\end{auditfinding}

\begin{auditfinding}[Observer-chain selection may leak preferred structure]
The observer chain \(\Gamma\) is defined as a coherently identifiable ordered
sequence under \(\PhiFlow\). This is plausible as a source object. The audit
risk is that the chain-selection rule may privilege the parameter \(\lambda\)
as an observable rest frame rather than a source ordering or gauge-like
bookkeeping variable. Any such leakage must either be operationally absent or
bounded as a testable deviation.
\end{auditfinding}

\section{Audit Classification}

\begin{repairburden}[Minimal repair burden]
A promotion-eligible successor packet must provide the following source-only
definitions:
\begin{enumerate}
    \item A front-generation rule for \(\Front\) from source transition
    structure.
    \item A coincidence-response protocol for \(\kappa_{\Gamma}\) that does
    not presuppose target spatial distance, simultaneity, radar coordinates,
    or target light-signal conventions.
    \item A clock-scale law for \(\alpha_{\Gamma}\) based on stable source
    cycles with explicit dimensional assignments.
    \item An invariance statement showing which relabelings of \(\Sub\),
    \(\lambda\), and response coordinates are redundant.
    \item A separation between source axioms and benchmark recovery tests.
\end{enumerate}
\end{repairburden}

\begin{auditverdict}[Promotion blocked, candidate not rejected]
The localization packet remains admissible as a draft candidate because it
explicitly marks the front, coincidence, and scale structures as source-side
and conditional. It is not admissible for promotion. The unresolved variables
\(\Front\), \(\kappa_{\Gamma}\), and \(\alpha_{\Gamma}\) are exactly the points
where target null structure and target clock structure could be smuggled into
the construction.
\end{auditverdict}

\section{Next Research Move}

The logical next step is not Gate Chair review. Two bounded paths are coherent:

\begin{enumerate}
    \item A Candidate Constructor repair job that supplies source-only
    definitions for \(\Front\), \(\kappa_{\Gamma}\), and \(\alpha_{\Gamma}\)
    under the blockers above.
    \item A Refuter job that assumes no repair and tests whether the current
    packet fails by unavoidable target import or underdetermination.
\end{enumerate}

The Director should choose between these paths in a separate continuation
invocation. The default conservative choice is repair before refutation if the
goal is constructive progress; the default adversarial choice is refutation if
the goal is early failure discovery.

\section*{References}

Ricciardi, A. S. (2026, June 8). \emph{Localization candidate packet: A
source-side observer-localization map for Gate 0 review} [Unpublished
manuscript]. The Aether-Flow Interpretation of Relativity Research Project,
\texttt{research\_control/tasks/RT-20260608-003/artifacts/02\_LOCALIZATION\_CANDIDATE\_PACKET.tex}.

Ricciardi, A. S. (2026, June 8). \emph{Gate 0 ontology formalization: Source
primitives, forbidden imports, and localization burden} [Unpublished
manuscript]. The Aether-Flow Interpretation of Relativity Research Project,
\texttt{research\_control/tasks/RT-20260608-002/artifacts/01\_GATE0\_ONTOLOGY\_FORMALIZATION.tex}.

\end{document}
